\subsection{The finite-support case and the complete finite-alteration classification}

\begin{statement}{A.69}{Every finite alteration of every integer at the actual level-three quotient}{Proved}
Let $E:y^2+y=x^3-1$, $R=(1,0)$, $U=E-E[3]$, and let
$\Gamma_3$ be the actual joint compact and modular j-line
quotient defined in A.67. Suppose $a\in\widehat\Z$ agrees
with an integer $n$ outside a finite set of coefficient primes.
Then $a\kappa(R)$ has a $\Gamma_3$ Selmer lift if and only if
$a=n$.

In particular, every nonzero exponent supported at finitely
many coefficient primes is excluded already by this actual
finite congruence refinement of the compact quotient. The zero
exponent is realized by the rational $O$-cusp. This completes
the finite-alteration exponent set, including the baseline
zero left open in A.68. It does not classify the finite
quotient's fibres or arbitrary infinite alterations.
\end{statement}
\begin{proof}
A.68 proves the assertion for $n\ne0$. It remains to exclude
every nonzero finite-support exponent. Put
\[
 K=\Q(\zeta),\quad \zeta^2+\zeta+1=0,\quad b^3=3,
 \quad L=K(b)=\Q(E[3]),\qquad
 g=\frac{b^2\zeta^2x-y-2}{b^2x-y-2}.
 \tag{A69.1}
\]
The quadratic Kummer character of $g$ factors through
$\Gamma_{3,L}$ by the direct kernel argument in A.67.
Its divisor is $3(d_1)-3(d_0)$, where
$d_i=(b\zeta^i,1)$, and it extends as a unit at $O$ with
$g(O)=1$. A hypothesized Selmer lift therefore supplies
one global class
$\beta\in H^1(L,\mu_2)=L^\times/L^{\times2}$.
At an interior local centre $Q$, its restriction is the
squareclass of $g(Q)$. At the centre $O$ it is trivial,
independently of the tangent parameter, because $g$ kills
the cusp inertia there.

\par\smallskip\noindent\textbf{Quadratic fields in a joint proper division tower.}\enspace
Let $\mathcal P$ be the nonempty finite support of $a$.
Take a finite abelian extension $F_0/K$ and form the
joint torsion and $R$-division tower $T_\infty/F_0$
for $p\in\mathcal P$. Its torsion subfield is
\[
 A_\infty=F_0(E[p^\infty]:p\in\mathcal P).
\]
The extension $A_\infty/K$ is abelian. Indeed all CM
endomorphisms are defined over $K$, and the Tate module
is a rank-one module over the CM algebra; its Galois
image is consequently commutative. The compositum
with the abelian field $F_0$ remains abelian.

The translation kernel is a closed additive subgroup
\[
 J=\operatorname{Gal}(T_\infty/A_\infty)
       \subset\prod_{p\in\mathcal P}T_pE.
 \tag{A69.2}
\]
As in A.61 it is the product of its primary subgroups
$J_p$, each stable under the joint torsion image.
Every quadratic character of
$\operatorname{Gal}(T_\infty/F_0)$ kills the odd
$J_p$. If $2\in\mathcal P$, suppose that the residual
$E[2]$ image over $F_0$ is $C_3$. Choose an element
$A$ of the joint torsion image whose two-adic
component has residual order three. Then
$\det(A-1)\in\Z_2^\times$. Cayley--Hamilton
expresses $(A-1)^{-1}$ as an integral polynomial
in $A$, so
\[
 (A-1)J_2=J_2.
 \tag{A69.3}
\]
A quadratic character kills the commutators
$[A,J_2]=(A-1)J_2$, and hence kills $J_2$ as well.
We have proved that every quadratic extension of
$F_0$ contained in $T_\infty$ lies in $A_\infty$.
In particular it is Galois, indeed abelian, over $K$.
When $2\notin\mathcal P$ the same conclusion holds
without a residual-image hypothesis.

Suppose also that $L\subset T_\infty$. Then
\[
 \sqrt\beta\in T_\infty,
 \tag{A69.4}
\]
where a representative of the squareclass is understood.
Otherwise, over the finite intermediate field
$F_0L$, the quadratic extension defined by $\beta$
is disjoint from the Galois tower. Apply the
simultaneous zero-primary condition of A.61 after
this base extension. Chebotarev permits that open
condition and nontrivial quadratic Frobenius
simultaneously. At the resulting good places,
$R_p=0$ for all $p\in\mathcal P$, so the actual
compact local point is $aR=O$. Full local Tate
Kummer injectivity identifies the Selmer centre
as $O$. But $g(O)=1$ and its cusp inertia is
killed, contradicting nontrivial quadratic
Frobenius. This proves (A69.4).

Consequently it suffices to find $p\in\mathcal P$
and a $p$-primary torsion point $S\notin E[3]$
such that
\[
 h=g(S)\in F_0^\times,
 \qquad F_0(\sqrt h)/K\ \hbox{is not Galois},
 \tag{A69.5}
\]
with the field and residual-image conditions above.
Indeed (A69.2)--(A69.3) then give
$\sqrt h\notin T_\infty$. Choose a $p$-primary
point $T$ with $a_pT=S$; this is possible by
dividing by the finite $p$-adic valuation of
$a_p\ne0$ and inverting its unit part. Take a
finite intermediate field $F_r\subset T_\infty$
containing $F_0L$, $T$, and $\sqrt\beta$.
Then $F_r(\sqrt{h/\beta})$ is quadratic and
disjoint from $T_\infty$. The prescribed-primary
lemma proved in A.68, reapplied over this field,
allows the simultaneous conditions
\[
 R_p=T,\qquad R_q=0\quad(q\in\mathcal P,\ q\ne p),
 \tag{A69.6}
\]
and nontrivial Frobenius on $\sqrt{h/\beta}$.
The forced compact point is $aR=S$, which is
interior. Local Selmer compatibility makes
$g(S)/\beta$ a square, a contradiction.

The order of these field arguments matters.
Noncontainment of $\sqrt h$ is proved over $F_0$
before any extension that kills the residual
$C_3$. It then persists under every finite
intermediate extension. Only (A69.4), not an
unproved assertion that $\beta$ remains in a
torsion subfield after base change, is needed.

\par\smallskip\noindent\textbf{An explicit coefficient-two witness.}\enspace
Suppose first that $2\in\mathcal P$. Set
$c^3=2$ and $F_0=K(c)$. The classes of $2$ and
$3$ are independent in $K^\times/K^{\times3}$:
valuation at the prime above two first forces
the exponent of $2$ to be divisible by three,
and $v_{1-\zeta}(3)=2$ then does the same for
the exponent of $3$. Thus $b\notin F_0$.
Put
\[
 r=b/c^2,\qquad r^3=3/4,
 \qquad S=(r,-1/2)\in E[2]\setminus\{O\}.
\]
We have $F_0(E[2])=F_0(b)=F_0L$. Therefore
$L\subset T_\infty$, and the residual
$E[2]$ image over $F_0$ is precisely $C_3$.
Direct substitution gives
\[
 h=g(S)=\frac{\zeta^2c-1}{c-1}\in F_0^\times.
 \tag{A69.7}
\]
At the prime defined by $c\equiv4$, $\zeta\equiv25$
modulo $31$, the two defining polynomials have
simple roots and
\[
 h\equiv27\pmod{31},\qquad 27^{15}\equiv-1\pmod{31}.
\]
Hence $h$ is a nonsquare in $F_0$.
Its three conjugates under $c\mapsto\zeta c$
have product one. If $F_0(\sqrt h)/K$ were
Galois, the squareclass of $h$ would be
invariant under this $C_3$ action. Their
product would then have squareclass
$[h]^3=[h]$, contradicting its being one.
This proves (A69.5), and excludes every
finite-support exponent with $a_2\ne0$.

\par\smallskip\noindent\textbf{Full Cartan orbits for every coefficient prime $p\ge5$.}\enspace
For such a prime the image on $E[p]$ over $K$
is the full Cartan $(\mathcal O_K/p)^\times$.
We give the local reason rather than assume
global CM surjectivity. At a split prime,
the two CM idempotents split the $p$-divisible
group into height-one factors, one formal
and one \'{e}tale. Ordinary inertia at one
prime above $p$ gives one full
$\mathbf F_p^\times$ factor. Inertia at the
conjugate prime gives the other, and together
they generate the full split Cartan.
At an inert prime the reduction is
supersingular: an ordinary connected
height-one factor could not carry the
unramified quadratic CM action.
The local extension $K_p/\Q_p$ is unramified,
so the $e=1$ supersingular inertia theorem
gives the full nonsplit Cartan
$\mathbf F_{p^2}^\times$.
These are the ordinary and supersingular
inertia statements in
\citet[Propositions 11--12, pp.~273--275]{AMSerre1972}.

In both cases the inertia groups at $p$
generate all of
$\operatorname{Gal}(K(E[p])/K)$.
The field $L/K$ is unramified at $p$, whence
\[
 K(E[p])\cap L=K,
 \qquad \operatorname{Gal}(L(E[p])/L)
           =(\mathcal O_K/p)^\times.
 \tag{A69.8}
\]
Choose $S\ne O$ if $p$ is inert, and choose
$S$ with both split CM components nonzero
if $p$ splits. Its orbit over $L$ has size
\[
 \begin{cases}
 p^2-1\ge24,&p\equiv2\pmod3,\\
 (p-1)^2\ge36,&p\equiv1\pmod3.
 \end{cases}
 \tag{A69.9}
\]
We next prove uniformly that $g(S)$ is a
nonsquare in $F_0=L(E[p])$.

\par\smallskip\noindent\textbf{A dyadic polynomial with only twenty-three zeros.}\enspace
The curve has good reduction at two and
$\#E(\mathbf F_2)=3$, so
$\operatorname{Frob}_4=[-2]$ on the reduction.
At a dyadic embedding choose the cube root
$b\in\Q_2$, for which $b\equiv3\pmod4$ and
$b^2\equiv1\pmod4$. The field $L$ splits
over $K$ at this place. Every completion of
$F_0$ above it is unramified over $\Q_2$:
$p$ is odd and the curve has good reduction.
On odd torsion the arithmetic Frobenius
$\sigma_4$ therefore acts as $[-2]$;
this follows from the same equality on
reduction and injectivity of odd-torsion
reduction.

If $g(A)=u^2$ for a nonzero $p$-torsion
point $A$ in this local field, then $u$
is a unit and $\sigma_4(u)\equiv u^4\pmod2$.
Squaring this congruence gives
\[
 g([-2]A)=\sigma_4(g(A))\equiv g(A)^4\pmod4.
 \tag{A69.10}
\]
All these values are units because their
zeros and poles are supported on $E[3]$,
which is disjoint from the reduction of
nonzero $p$-torsion.

Replace $b^2$ by one modulo four and write
\[
 g_0=\frac{\zeta^2x-y-2}{x-y-2},\qquad
 m=\frac{3x^2}{2y+1},\qquad
 X=m^2-2x,\quad Y=y-m(x-X).
\]
Here $(X,Y)=[-2](x,y)$. The exact numerator
of $g_0(X,Y)-g_0(x,y)^4$, reduced using
$y^2+y=x^3-1$ and $\zeta^2+\zeta+1=0$,
is divisible by two. Its half modulo two is
\begin{align*}
 F={}&(\zeta+1)x^{11}+\zeta x^{10}y+x^9
       +(\zeta+1)x^8y+x^6y+x^6\\
    &\quad+\zeta x^4y+(\zeta+1)x^2+\zeta xy.
 \tag{A69.11}
\end{align*}
The corresponding reduced denominator
modulo two is $(x+y)^8$.
On the affine curve in characteristic two,
$x+y=0$ forces $(x+1)^3=0$; its only
geometric zero is the order-three point
$(1,1)$. Thus the denominator is a unit
at all the torsion points in question,
and (A69.10) implies $F(\overline A)=0$.

The polynomial $F$ is a nonzero function
on $E/\mathbf F_4$ with exact pole order
$23$ at $O$. Indeed $x$ has pole order
two, $y$ has pole order three, and
$\zeta x^{10}y$ is the unique term of
pole order $23$. There are no other poles,
so $F$ has at most $23$ geometric zeros.
If $g(S)$ were a square in $F_0$, all its
conjugate values under the Cartan group
in (A69.8) would also be squares.
Their reductions at the chosen dyadic
place are distinct, since reduction is
injective on odd torsion. Formula
(A69.9) would therefore give at least
$24$ distinct zeros of $F$, a contradiction.
This proves the nonsquare assertion
for every $p\ge5$, with no numerical
range restriction.

\par\smallskip\noindent\textbf{The cubic conjugation and conclusion.}\enspace
The field $F_0=L(E[p])$ is abelian over
$K$. By (A69.8) it has an order-three
automorphism $\tau$ fixing $S$ and
sending $b$ to $\zeta b$. The three
tangent-line ratios telescope, giving
\[
 g(S)\,\tau(g(S))\,\tau^2(g(S))=1.
 \tag{A69.12}
\]
Since $h=g(S)$ is nonsquare, the same
squareclass argument as for (A69.7)
shows that $F_0(\sqrt h)/K$ is not
Galois. Moreover $F_0$ is unramified
at two, whereas a nonzero point of
$E[2]$ requires a cube root of $3/4$,
of dyadic valuation $-2/3$. The
residual $E[2]$ image over $F_0$
therefore remains $C_3$. All the
hypotheses of (A69.5)--(A69.6) hold.
This excludes every finite-support
exponent having a nonzero component
at some $p\ge5$, even if other
components are present.

Every nonzero finite-support exponent
has a nonzero two-component, a nonzero
$p\ge5$ component, or support exactly
$\{3\}$. The first two cases were
excluded above and the third in A.68.
The zero exponent has the rational
$O$-cusp lift. For nonzero integer
$n$, the point $nR$ is interior and
provides the required lift. This
completes both directions.

The exact symbolic identities,
denominator reduction, pole calculation,
cubic norm identities and nonsquare
certificate at $31$ are checked in the
executed program
\texttt{code/a\_main\_finite\_support.py};
its raw output is saved alongside it.
The calculation is a certificate for
the displayed uniform proof, not a
search through a bounded list of
coefficient primes.
\end{proof}
